% -------------------------------------------------
% Anexo I. Resultados completos de las encuestas aplicadas
% -------------------------------------------------
\section*{Anexo I. Resultados completos de las encuestas aplicadas}
\addcontentsline{toc}{section}{Anexo I. Resultados completos de las encuestas aplicadas}
\label{anexo:encuestas}

El presente anexo reúne la evidencia completa correspondiente a las tres encuestas aplicadas durante la fase de evaluación del prototipo.
\\
\\
Con el fin de mantener el capítulo de resultados enfocado en el análisis, aquí se incluyen los documentos originales generados por Google Forms, así como las tablas completas con los resultados por ítem, promedios y número de respuestas. Cada encuesta aparece organizada según su propósito y población objetivo.

\vspace{0.4cm}

% -------------------------------------------------
% Encuesta 1
% -------------------------------------------------
\subsection*{Encuesta 1: Evaluación del módulo \textit{Aprender a Aprender}}
\addcontentsline{toc}{subsection}{Encuesta 1: Evaluación del módulo Aprender a Aprender}

\vspace{0.2cm}
\noindent\textbf{Archivo anexo:}  
\texttt{Resultados Encuesta 1 - Modulo Aprender a Aprender.pdf}

\vspace{0.5cm}

% ---------- Tablas de la Encuesta 1 ----------
\subsubsection*{Tablas de resultados por ítem}
\addcontentsline{toc}{subsubsection}{Encuesta 1: Tablas de resultados}

% --- Comprensión y uso ---
\begin{table}[H]
\centering
\footnotesize
\setlength{\tabcolsep}{4pt}
\renewcommand{\arraystretch}{1.35}
\begin{tabular}{|p{10.5cm}|c|c|}
\hline
\textbf{Ítem (pregunta literal)} & \textbf{Promedio} & \textbf{N} \\
\hline
Comprendí claramente cómo funcionan las metas de aprendizaje, los objetivos, las tareas y las rutas o planes de aprendizaje dentro del módulo. & 4.32 & 38 \\
\hline
Entendí cómo utilizar el tablero Kanban para organizar y visualizar mis tareas. & 4.26 & 38 \\
\hline
Comprendí cómo usar el chatbot para generar un roadmap de aprendizaje personalizado. & 4.00 & 38 \\
\hline
Entendí el propósito de las autoevaluaciones y reportes de progreso. & 4.39 & 38 \\
\hline
\end{tabular}
\caption{Encuesta 1 – Comprensión y uso de las herramientas del módulo.}
\label{tab:anexo-aaa-comprension}
\end{table}

% --- Motivación y participación ---
\begin{table}[H]
\centering
\footnotesize
\setlength{\tabcolsep}{4pt}
\renewcommand{\arraystretch}{1.35}
\begin{tabular}{|p{10.5cm}|c|c|}
\hline
\textbf{Ítem (pregunta literal)} & \textbf{Promedio} & \textbf{N} \\
\hline
Los procesos de crear metas, objetivos y tareas me mantuvieron participando activamente. & 4.29 & 38 \\
\hline
El tablero Kanban me ayudó a mantenerme activo/a moviendo tareas entre estados. & 4.13 & 38 \\
\hline
El editor visual de roadmaps me permitió involucrarme de forma dinámica con mi plan de aprendizaje. & 4.34 & 38 \\
\hline
La creación de roadmaps con el chatbot fue una experiencia interactiva que me mantuvo comprometido/a. & 4.21 & 38 \\
\hline
El módulo me motivó a trabajar activamente en mis metas de aprendizaje. & 4.29 & 38 \\
\hline
Tener seguimiento visual del progreso (dashboard, porcentajes, estados de tareas) aumentó mi motivación. & 4.42 & 38 \\
\hline
Poder personalizar mis tareas (tipos, recursos, notas) incrementó mi interés en continuar. & 4.39 & 38 \\
\hline
\end{tabular}
\caption{Encuesta 1 – Motivación y participación.}
\label{tab:anexo-aaa-motivacion}
\end{table}

% --- Aplicación práctica ---
\begin{table}[H]
\centering
\footnotesize
\setlength{\tabcolsep}{4pt}
\renewcommand{\arraystretch}{1.35}
\begin{tabular}{|p{10.5cm}|c|c|}
\hline
\textbf{Ítem (pregunta literal)} & \textbf{Promedio} & \textbf{N} \\
\hline
Siento que retuve los conceptos clave sobre planificación, autorregulación y organización del aprendizaje. & 4.32 & 37 \\
\hline
Las autoevaluaciones me ayudaron a recordar mejor mi progreso y aprendizajes. & 4.08 & 37 \\
\hline
El roadmap generado me facilitó recordar los pasos de mi plan de estudio. & 4.38 & 37 \\
\hline
Puedo aplicar en la vida real los conceptos del módulo. & 4.54 & 37 \\
\hline
El roadmap generado me sirve como guía práctica para continuar mi aprendizaje. & 4.27 & 37 \\
\hline
\end{tabular}
\caption{Encuesta 1 – Retención y aplicación práctica.}
\label{tab:anexo-aaa-aplicacion}
\end{table}

% -------------------------------------------------
% Encuesta 2
% -------------------------------------------------
\subsection*{Encuesta 2: Evaluación del módulo \textit{Escucha Activa}}
\addcontentsline{toc}{subsection}{Encuesta 2: Evaluación del módulo Escucha Activa}

\vspace{0.2cm}
\noindent\textbf{Archivo anexo:}  
\texttt{Resultados Encuesta 2 - Modulo Escucha Activa.pdf}

\vspace{0.5cm}

\subsubsection*{Tablas de resultados por ítem}
\addcontentsline{toc}{subsubsection}{Encuesta 2: Tablas de resultados}

% --- Comprensión ---
\begin{table}[H]
\centering
\footnotesize
\setlength{\tabcolsep}{4pt}
\renewcommand{\arraystretch}{1.35}
\begin{tabular}{|p{10.5cm}|c|c|}
\hline
El módulo me ayudó a comprender los principios de la escucha activa y cómo se aplican a situaciones reales. & 4.80 & 25 \\
\hline
Las explicaciones y retroalimentación después de cada ronda fueron claras. & 4.64 & 25 \\
\hline
\end{tabular}
\caption{Encuesta 2 – Comprensión del contenido.}
\label{tab:anexo-ea-comprension}
\end{table}

% --- Motivación ---
\begin{table}[H]
\centering
\footnotesize
\setlength{\tabcolsep}{4pt}
\renewcommand{\arraystretch}{1.35}
\begin{tabular}{|p{10.5cm}|c|c|}
\hline
Las actividades del juego fueron dinámicas y me mantuvieron participando. & 4.72 & 25 \\
\hline
Los distintos modos de juego hicieron la experiencia variada. & 4.48 & 25 \\
\hline
La estructura por rondas, puntajes y feedback me motivó. & 4.52 & 25 \\
\hline
La opción de replay me motivó a mejorar. & 4.44 & 25 \\
\hline
\end{tabular}
\caption{Encuesta 2 – Motivación y participación.}
\label{tab:anexo-ea-motivacion}
\end{table}

% --- Retención ---
\begin{table}[H]
\centering
\footnotesize
\setlength{\tabcolsep}{4pt}
\renewcommand{\arraystretch}{1.35}
\begin{tabular}{|p{10.5cm}|c|c|}
\hline
Siento que mejoré mi capacidad para recordar detalles importantes. & 92\% Sí, 8\% No & 25 \\
\hline
\end{tabular}
\caption{Encuesta 2 – Retención de información (dicotómica).}
\label{tab:anexo-ea-retencion-dicotomica}
\end{table}

\begin{table}[H]
\centering
\footnotesize
\setlength{\tabcolsep}{4pt}
\renewcommand{\arraystretch}{1.35}
\begin{tabular}{|p{10.5cm}|c|c|}
\hline
El uso de diferentes estilos de contenido fortaleció mi comprensión auditiva. & 4.52 & 25 \\
\hline
\end{tabular}
\caption{Encuesta 2 – Retención (Likert).}
\label{tab:anexo-ea-retencion-likert}
\end{table}

% --- Aplicación ---
\begin{table}[H]
\centering
\footnotesize
\setlength{\tabcolsep}{4pt}
\renewcommand{\arraystretch}{1.35}
\begin{tabular}{|p{10.5cm}|c|c|}
\hline
Puedo aplicar técnicas de escucha activa en situaciones reales. & 4.72 & 25 \\
\hline
El proceso del módulo me ayudó a desarrollar habilidades transferibles. & 4.44 & 25 \\
\hline
\end{tabular}
\caption{Encuesta 2 – Aplicación práctica.}
\label{tab:anexo-ea-aplicacion}
\end{table}

% --- Evaluación global ---
\begin{table}[H]
\centering
\footnotesize
\setlength{\tabcolsep}{4pt}
\renewcommand{\arraystretch}{1.35}
\begin{tabular}{|p{10.5cm}|c|c|}
\hline
La experiencia general fue útil, clara y bien diseñada. & 4.84 & 25 \\
\hline
Recomendaría la aplicación a otros estudiantes. & 4.92 & 25 \\
\hline
\end{tabular}
\caption{Encuesta 2 – Evaluación global.}
\label{tab:anexo-ea-global}
\end{table}
